	\chapter{Spectra}

In this chapter the notion of spectra is introduced. This is the set of prime ideals of a ring which naturally can be equipped with a topology, called the Zariski topology. We discuss spectra of familiar rings and the connection of this definition to classical algebraic geometry. Further topological properties of spectra are discussed; the most important one is the fact that the spectrum of a ring is quasi-compact. Finally we prove that mapping a ring to its spectrum defines a functor from the category of rings to the category of topological spaces.

Recall our convention that rings are required to be commutative and unital, and every ring homomorphism is unital.

\section{Spectrum of a Ring}
\begin{definition}
	Given a ring $R$, its \inlinedef{(prime) spectrum} $\Spec R$ is the set of its prime ideals.
\end{definition}

It turns out that the spectrum $\Spec R$ comes with a natural topology. To define this topology, we introduce the following notation. Given a subset $S \subseteq R$, define
\begin{equation*}
	V(S) \coloneqq \Set{ \mathfrak{p} \in \Spec R : \mathfrak{p} \supseteq S}.
\end{equation*}
It is straightforward to check that if $\mathfrak{a}$ is the ideal generated by $S$, then $V(S) = V(\mathfrak{a})$. Hence it suffices to consider only the sets $V(\mathfrak{a})$ whenever $\mathfrak{a}$ is an ideal of $R$. We shall also write $V(f_1, \dotsc, f_n)$ for $V(\Gen{f_1, \dotsc, f_n})$, where, as per usual, $\Gen{f_1, \dotsc, f_n}$ denotes the ideal generated by $f_1, \dotsc, f_n$. 

\begin{exercise}
	Show that for any ideal $\mathfrak{a}$ of $R$, $V(\sqrt{\mathfrak{a}}) = V(\mathfrak{a})$, where $\sqrt{\mathfrak{a}}$ denotes the radical of $\mathfrak{a}$.
\end{exercise}

The sets $V(\mathfrak{a})$ satisfy the following properties.

\begin{proposition}\label{chap2:zariski-topology}
	Let $R$ be a ring, and let $\mathfrak{a}$, $\mathfrak{b}$ be ideals of $R$. Further, for each $k \in S$ let $\mathfrak{a}_k$ be ideals of $R$ where $S$ is some index set. Then
	
	\begin{enumerate}
		\item If $\mathfrak{a} \subseteq \mathfrak{b}$, then $V(\mathfrak{b}) \subseteq V(\mathfrak{a})$.
		
		\item $V(\emptyset) = \Spec R$ and $V(1) = \emptyset$.
		
		\item $V(\mathfrak{a} \cap \mathfrak{b}) = V(\mathfrak{a}\mathfrak{b}) = V(\mathfrak{a}) \cup V(\mathfrak{b})$.
		
		\item $V(\sum_{k \in S} \mathfrak{a}_k) = \bigcap_{k \in S} V(\mathfrak{a}_k)$.
	\end{enumerate}
\end{proposition}
\begin{exercise}
	Prove the above proposition. \textit{[Hint: the fact that $\sqrt{\mathfrak{a}\mathfrak{b}} = \sqrt{\mathfrak{a} \cap \mathfrak{b}}$ might be useful.]}
\end{exercise}

The above proposition shows that the sets of the form $V(\mathfrak{a})$, where $\mathfrak{a}$ runs through all ideals of $R$, are closed sets of a topology on $\Spec R$. This is the \inlinedef{Zariski topology} on $\Spec R$. With the topology in mind, we will refer to elements of $\Spec R$ as points. Given $x \in \Spec R$, we shall write $\mathfrak{p}_x$ if we want to view $x$ as a prime ideal of $R$.

There is a reverse correspondence of the map $V(-)$. If $Y$ is a subset of $\Spec R$, we define the ideal $I(Y)$ associated to $Y$ by
\begin{equation*}
	I(Y) \coloneqq \bigcap_{\mathfrak{p} \in Y} \mathfrak{p}.
\end{equation*}
Note that $I(\emptyset) = R$. 

\begin{proposition}\label{chap2:I-and-V-correspondence}
	Let $R$ be a ring, $\mathfrak{a} \subseteq R$ an ideal, and $Y$ a subset of $\Spec R$.
	\begin{enumerate}
		\item If $X \subseteq Y \subseteq \Spec R$, then $I(Y) \subseteq I(X)$.
		\item $\sqrt{I(Y)} = I(Y)$, i.e. $I(Y)$ is a radical ideal.
		\item $I(V(\mathfrak{a})) = \sqrt{\mathfrak{a}}$ and $V(I(Y)) = \overline{Y}$ where $\overline{Y}$ is the topological closure of $Y$ in $\Spec R$.
	\end{enumerate}
\end{proposition}

\begin{proof}
	(1.) Trivial.
	
	(2.) We only need to check the non-trivial inclusion $\sqrt{I(Y)} \subseteq I(Y)$. Suppose $a \in R$ and $n \geq 1$ satisfy $a^n \in I(Y)$. Then for each prime ideal $\mathfrak{p} \in Y$, $a^n \in \mathfrak{p}$. Then $a \in \mathfrak{p}$, and so $a \in I(Y)$.
	
	(3.) The assertion that $I(V(\mathfrak{a})) = \sqrt{\mathfrak{a}}$ follows from a theorem in commutative algebra which states that, for any ideal $\mathfrak{a}$ of $R$,
	\begin{equation*}
		\sqrt{\mathfrak{a}} = \bigcap_{\substack{\mathfrak{p} \text{ prime} \\ \mathfrak{p} \supseteq \mathfrak{a}}} \mathfrak{p}.
	\end{equation*}
	For the second assertion, note that $V(I(Y))$ contains $Y$. Now if $V(\mathfrak{b})$ is a closed set containing $Y$, then $I(Y) \supseteq I(V(\mathfrak{b})) = \sqrt{\mathfrak{b}} \supseteq \mathfrak{b}$ by the first assertion. It follows that 
	\begin{equation*}
		V(I(Y)) \subseteq V(\mathfrak{b}),
	\end{equation*}
	proving that $V(I(Y))$ is the closure of $Y$.
\end{proof}

The following corollaries are immediate.
\begin{corollary}\label{chap2:V-and-I-inverse}
	The maps $Y \mapsto I(Y)$ and $\mathfrak{a} \mapsto V(\mathfrak{a})$ give rise to mutually inverse and inclusion-reversing bijections between the sets
	\begin{equation*}
		\Set*{\text{radical ideals of $R$}} \rightleftarrows \Set*{\text{closed subsets of $\Spec R$}}.
	\end{equation*}
\end{corollary}

\begin{corollary}
	Let $R$ be a ring and $\mathfrak{a}, \mathfrak{b}$ ideals of $R$. Then $V(\mathfrak{a}) = V(\mathfrak{b})$ if and only if $\sqrt{\mathfrak{a}} = \sqrt{\mathfrak{b}}$.
\end{corollary}

%	We also draw the following consequence of. The singleton $\Set{\mathfrak{p}_x} \subseteq \Spec R$ is a closed subset if and only if $V(\mathfrak{p}_x) = \Set{\mathfrak{p}_x}$, if and only if $\mathfrak{p}_x$ is a maximal ideal.

There is a connection between the sets $V(I)$ with the new definition in this chapter and the vanishing sets in the sense of classical algebraic geometry. For this, elements of $R$ need to be interpreted as ``functions'' on $\Spec R$. Given $f \in R$, we define \inlinedef{the value of $f$ at a prime} $\mathfrak{p} \in \Spec R$ to be the image of $f$ under the composition
\begin{equation*}
	R \to R/\mathfrak{p} \hookrightarrow \Frac(R/\mathfrak{p}).
\end{equation*}
Here we used the fact that $R/\mathfrak{p}$ is an integral domain so that its field of fractions may be constructed. Via this map, we see that $f(\mathfrak{p}) = 0$ if and only if $f \in \mathfrak{p}$. Hence $V(f)$ consists of prime ideals $\mathfrak{p}$ such that $f(\mathfrak{p}) = 0$. Therefore $V(f)$ is the \textit{vanishing set} of $f$. Compare this situation to the classical case when $R$ is the polynomial ring $k[x_1, \dotsc, x_n]$ and $f$ is a polynomial in $R$.

Here we discuss spectra of familiar rings.

\begin{example}[Spectrum of a Field]
	Every field $k$ has only one prime ideal $\Set{0}$. Thus $\Spec k$ consists of a single point.
\end{example}

\begin{example}[Spectrum of $\ZZ$]
	The prime ideal of $\ZZ$ are the zero ideal $(0)$ and the ideal generated by prime numbers $(p)$. Hence $\ZZ$ has infinitely many points.
\end{example}

More generally, the spectrum of an integral domain $R$ is non-empty. It contains the point corresponding to the zero ideal. Let us denote this point by $\eta$. The fact that every prime ideal of $R$ contains the zero ideal means that $\overline{\Set{\eta}} = \Spec R$: the closure of the singleton $\Set{\eta}$ is equal to the whole spectrum. Such a point is called a \textit{generic point} of $\Spec R$ (\cref{chap2:spacial-points-of-rings}). As strange as it is, generic points will turn out to be important to the theory of integral schemes.





For the last example in this section, we relate the coordinate ring $k[V]$ of an affine variety $V \subseteq \AAA^n$ to the spectrum $\Spec k[V]$ when $k$ is an algebraically closed field.

\begin{example}
	Let $V \subseteq \AAA^n$ be an affine variety given as the vanishing set $Z(f_1, \dotsc, f_i)$ of polynomials in $k[x_1, \dotsc, x_n]$. Let $I$ denote the ideal consisting of polynomials vanishing on $V$. Recall that there is a one-to-one correspondence
	\begin{equation*}
		\Set*{\text{maximal ideals of $k[x_1, \dotsc, x_n]/I$}} \Longleftrightarrow \Set*{\text{points of $V$}}.
	\end{equation*}
	Let us denote the coordinate ring of $V$ by $k[V] =k[x_1, \dotsc, x_n]/I$. If we write $\Spec_{M}(A)$ for the set of maximal ideals of the ring $A$, then the points of $V$ are in bijection with points of $\Spec_{M}(k[V])$ via Hilbert Nullstellensatz. The set $\Spec_{M}(k[V])$ is a subspace of $\Spec(k[V])$. It is easy to check that the subspace topology on $\Spec_{M}(k[V])$ is the same one as the Zariski topology on $V$ considered as a subspace of $k[x_1, \dotsc, x_n]$.
\end{example}

In topological terms maximal ideals are the closed points of the spectrum (see \cref{chap2:spacial-points-of-rings} below). Hence we ``recover'' classical algebraic geometry when we consider the closed points of the spectrum of a coordinate ring.

\section{Topological Properties of Spectra}

\subsection*{Basis of $\Spec R$}
We now turn to the topological properties of $\Spec R$ under the Zariski topology. Firstly, there is a basis for this topology which is derived from elements of $R$.

\begin{definition}
	Let $R$ be a ring and $f \in R$. We define
	\begin{equation*}
		D(f) = \Spec R \setminus V(f).
	\end{equation*} 
	Plainly, $D(f)$ is open in $\Spec R$. An open set of this form is called a \inlinedef{principal open subset} of $\Spec R$.
\end{definition}

The set $D(f)$ can also be regarded as the set of points where $f$ does not vanish, that is, $D(f) = \Set{\mathfrak{p} \in \Spec R : f(\mathfrak{p}) \neq 0}$.

\begin{proposition}
	Let $R$ be a ring.	Principal open subsets of $\Spec R$ form a basis of the Zariski topology.
\end{proposition}
\begin{proof}
	Let $U \subseteq \Spec R$ be an open set, so that $U = \Spec R \setminus V(\mathfrak{b})$ for some ideal $\mathfrak{b}$ of $R$. It is easy to see that
	\begin{equation*}
		V(\mathfrak{b}) = \bigcap_{f \in \mathfrak{b}} V(f),
	\end{equation*}
	and consequently $U = \bigcup_{f \in \mathfrak{b}} D(f)$.
\end{proof}

The basis $\mathcal{B}$ consisting of principal open subsets is \textit{stable under finite intersections}. This means that if $B_1, B_2 \in \mathcal{B}$, then $B_1 \cap B_2 \in \mathcal{B}$. This follows from the proposition below, and the fact that $D(0) = \emptyset$.

\begin{proposition}\label{chap2:principal-open-is-basis}
	Let $R$ be a ring. Then $D(f) \cap D(g) = D(fg)$ for all $f, g\in R$. Consequently the basis consisting of principal open subsets is stable under finite intersections.
\end{proposition}
\begin{exercise}
	Prove the above proposition.
\end{exercise}

The following lemma is useful for turning a topological statement about a covering by principal open subsets into an algebraic statement.

\begin{lemma}[``Covering Trick'']\label{chap2:covering-trick}
	Let $\Set{f_i}_{i \in I}$ be a family of elements of a ring $R$ and let $g \in R$. Then $D(g) \subseteq \bigcup_{i \in I}D(f_i)$ if and only if there is an integer $n >0$ such that $g^n$ is contained in the ideal $\mathfrak{a}$ generated by $f_i$.
\end{lemma}
\begin{proof}
	The condition $D(g) \subseteq \bigcup_{i \in I}D(f_i)$ is equivalent to $V(g) \supseteq \bigcap_{i \in I} V(f_i) = V(\mathfrak{a})$, which is equivalent to $g \in \sqrt{\mathfrak{a}}$.
\end{proof}

Two special cases of the above theorem are often used. Firstly, $D(g) \subseteq D(f)$ if and only if $g \in \sqrt{(f)}$, if and only if $g^n = fa$ for some $a \in R$. This is important when we define the structure sheaf of $\Spec R$. 

Secondly, setting $g = 1$ shows that the collection $\Set{D(f_i)}_{i \in I}$ covers $\Spec R$ if and only if $1 \in \Gen{f_i : i \in I}$. Hence there is a finite subset $J \subseteq I$ such that $1 = \sum_{j \in J} r_j f_j$ where $r_j \in R$. This implies that $\Spec R = \bigcup_{j \in J} D(f_j)$. Thus we have found a finite subcover of $\Set{D(f_i)}_{i \in I}$, implying that $\Spec R$ is quasi-compact. More generally, every principal open subset of $\Spec R$ is quasi-compact by the same argument.

\begin{proposition}\label{chap2:principal-open-is-quasi-compact}
	Let $R$ be a ring. Then every principal open subset $D(g)$, $g \in R$, is quasi-compact. 
\end{proposition}
\begin{exercise}
	Prove this proposition by adapting the proof that $\Spec R$ is quasi-compact. 
\end{exercise}

\subsection*{Irreducibility}

A non-empty topological space is \inlinedef{irreducible} if it cannot be written as a union of two proper closed subsets. In other words, a topological space $X \neq \emptyset$ is irreducible if and only if whenever $X = C_1 \cup C_2$, where $C_1$ and $C_2$ are closed subsets, then either $C_1 = X$ or $C_2 = X$. A topological space is \inlinedef{reducible} if it is not irreducible. It is useful to consider the empty topological space as being reducible.

A subset of a topological space is irreducible if it is irreducible under the subspace topology. Similar to connected components, we define an \inlinedef{irreducible component} to be a subspace $Z \subseteq X$ such that $Z$ is maximally irreducible subset of $X$, that is, $Z$ is irreducible and whenever $Y$ is irreducible and $Z \subseteq Y$, then $Z = Y$.

\begin{exercise}
	Show that if $X$ is a topological space and $x \in X$, then the closure $\overline{\Set{x}}$ is irreducible.
\end{exercise}

We collect some properties of irreducible spaces here, whose proofs are mechanical.

\begin{proposition}
	The following are equivalent for a non-empty topological space $X$.
	\begin{enumerate}
		\item $X$ is irreducible.
		
		\item Every non-empty open subset of $X$ is dense in $X$.
		
		\item Every two non-empty open subsets of $X$ has non-empty intersection.
		
		\item Every non-empty open subset of $X$ is irreducible.
	\end{enumerate}
\end{proposition}

\begin{exercise}
	Prove the above proposition.
\end{exercise}

\begin{proposition}
	Let $f\colon X \to Y$ be a continuous map between topological spaces. If $Z \subseteq X$ is an irreducible subset, then its image $f(Z)$ is irreducible.
\end{proposition}

\begin{proof}
	Clearly $f(Z)$ is non-empty because $Z$ is. Suppose $f(Z) = C_1 \cup C_2$ where $C_1$ and $C_2$ are closed subsets of $f(Z)$ under the subspace topology. Hence $C_1 = D_1 \cap f(Z)$ and $C_2 = D_2 \cap f(Z)$ for some closed sets $D_1$ and $D_2$ of $Y$. We then have
	\begin{equation*}
		Z = (f^{-1}(D_1) \cap Z) \cup (f^{-1}(D_2) \cap Z).
	\end{equation*}
	This is a decomposition of $Z$ into two closed subsets. Since $Z$ is irreducible, either $f^{-1}(D_1) \cap Z = Z$ or $f^{-1}(D_2) \cap Z = Z$. Hence either $C_1 = f(Z)$ or $C_2 = f(Z)$.		
\end{proof}

\begin{proposition}\label{chap2:space-irreducible-iff-closure-irreducible}
	Let $X$ be a topological space. Then $Y\subseteq X$ is irreducible if and only if its closure $\overline{Y}$ is irreducible.
\end{proposition}

\begin{proof}
	First assume that $Y \subseteq X$ is irreducible. Then $\overline{Y}$ is clearly non-empty. If $\overline{Y} = C_1 \cup C_2$ for some closed subsets $C_1, C_2$ of $X$, then $Y = (C_1 \cap Y) \cup (C_2 \cap Y)$. Since $Y$ is irreducible, one of the subsets must be equal to $Y$, say $C_1 \cap Y = Y$. Hence $C_1$ is a closed set containing $Y$, and so $\overline{Y} \subseteq C_1$. Plainly, $C_1 \subseteq \overline{Y}$ and therefore $\overline{Y} = C_1$.
	
	For the converse, assume that $Y$ is reducible. If $Y$ is empty, then $\overline{Y} = \emptyset$ is also reducible. Hence we may assume that $Y$ is non-empty and write $Y = C_1 \cup C_2$, where $C_1$ and $C_2$ are proper closed subsets of $Y$ under the subspace topology. Then $\overline{Y} = \overline{C_1} \cup \overline{C_2}$ by the property of closure. If $\overline{Y} = \overline{C_1}$, then 
	\begin{equation*}
		C_1 = \overline{C_1} \cap Y = \overline{Y} \cap Y = Y,
	\end{equation*} 
	a contradiction. Similarly we conclude that $\overline{Y} \neq \overline{C_2}$. Therefore we have found a decomposition of $\overline{Y}$ into two proper closed subsets.
\end{proof}

We now turn to irreducible subsets of $\Spec R$, when $R$ is a ring. The starting point is the following proposition which says that irreducible subsets are related to prime ideals.

\begin{proposition}\label{chap2:irreducible-iff-prime}
	A subset $Y \subseteq \Spec R$ is irreducible if and only if $I(Y)$ is a prime ideal of $R$.
\end{proposition}

\begin{proof}
	Assume that $Y$ is irreducible. Then $\overline{Y}$ is also irreducible. Let $fg \in I(Y) \subseteq R$. Using \cref{chap2:zariski-topology} and \cref{chap2:I-and-V-correspondence}, we have the identities
	\begin{equation*}
		V(f) \cup V(g) = V(fg) \supseteq V(I(Y)) = \overline{Y}.
	\end{equation*} 
	Hence $\overline{Y} = (V(f) \cap \overline{Y}) \cup (V(g) \cap \overline{Y})$. Since $\overline{Y}$ is irreducible, it must be equal to one of the subsets, say $\overline{Y} = V(f) \cap \overline{Y}$. Therefore $\overline{Y} \subseteq V(f)$. Again, by \cref{chap2:I-and-V-correspondence}, we have $I(Y) \supseteq \sqrt{(f)}$. This implies that $f \in I(Y)$. This proves that $I(Y)$ is prime.
	
	Now assume that $\mathfrak{p} = I(Y)$ is prime. By \cref{chap2:I-and-V-correspondence}, $\overline{Y} = V(\mathfrak{p})$. Let $V(\mathfrak{a}), V(\mathfrak{b})$ be closed sets such that
	\begin{equation*}
		V(\mathfrak{p}) = V(\mathfrak{a}) \cup V(\mathfrak{b}) = V(\mathfrak{a}\cap \mathfrak{b}).
	\end{equation*}
	Recall that  $\mathfrak{p} = \sqrt{\mathfrak{a} \cap \mathfrak{b}} = \sqrt{\mathfrak{a}} \cap \sqrt{\mathfrak{b}}$ by commutative algebra. Since $\mathfrak{p}$ is prime, either $\mathfrak{p} = \sqrt{\mathfrak{a}}$ or $\mathfrak{p} = \sqrt{\mathfrak{b}}$. Hence $V(\mathfrak{p}) = V(\mathfrak{a})$ or $V(\mathfrak{p}) = V(\mathfrak{b})$. Thus $V(\mathfrak{p}) = \overline{Y}$ is irreducible, and so does $Y$.
\end{proof}

\begin{proposition}\label{chap2:correspondence-closed-irreducible-minimal}
	There is a one-to-one correspondence between $\Spec R$ and the set of closed irreducible subsets of $\Spec R$ given in one direction by $\mathfrak{p} \mapsto \overline{\Set{\mathfrak{p}}}$. Under this correspondence, minimal prime ideals of $R$ correspond bijectively to irreducible components of $\Spec R$.
\end{proposition}
\begin{proof}
	If $\mathfrak{p}$ is prime, then $I(\Set{\mathfrak{p}}) = \mathfrak{p}$ and thus $\overline{\mathfrak{\Set{p}}} = V(I(\Set{\mathfrak{p}})) = V(\mathfrak{p})$ is irreducible by \cref{chap2:irreducible-iff-prime}. Now if $Y$ is a closed irreducible subset of $\Spec R$, then by \cref{chap2:irreducible-iff-prime}, $I(Y)$ is a prime ideal. Hence the correspondence is given by the functions $V(-)$ and $I(-)$ which are proved to be mutually inverse bijections in \cref{chap2:V-and-I-inverse}. 
	
	The second statement follows immediately from the observation that $\overline{\Set{\mathfrak{p}}} \subseteq \overline{\Set{\mathfrak{q}}}$ if and only if $\mathfrak{p} \supseteq \mathfrak{q}$ for all prime ideals $\mathfrak{p}, \mathfrak{q}$ of $R$. 
\end{proof}

\subsection*{Special Points}

We introduce more definitions for points of a topological space which has interesting properties.

\begin{definition}
	Let $X$ be a topological space.
	\begin{enumerate}
		\item A point $x \in X$ is a \inlinedef{closed point} if $\Set{x}$ is closed.
		\item A point $\eta \in X$ is a \inlinedef{generic point} if $\overline{\Set{\eta} }= X$.
		
		\item Let $x$ and $x'$ be points of $X$. We say that $x$ is a \inlinedef{generization} of $x'$ and that $x'$ is a \inlinedef{specialisation} of $x$ if $x' \in \overline{\Set{x}}$.
		
		\item A point $x \in X$ is a \inlinedef{maximal point} of $X$ if $\overline{\Set{x}}$ is an irreducible component of $X$.
	\end{enumerate}
\end{definition}

For spectra of rings, the topological notions above correspond to the following algebraic statements.

\begin{proposition}\label{chap2:spacial-points-of-rings} Let $R$ be a ring.
	\begin{enumerate}
		\item A point $x \in \Spec R$ is closed if and only if the ideal $\mathfrak{p}_x$ is maximal.
		
		\item A point $\eta \in \Spec R$ is generic if and only if the ideal $\mathfrak{p}_{\eta}$ is the unique minimal prime ideal.
		
		\item A point $x$ is a generization of $x'$ if and only if $\mathfrak{p}_x \subseteq \mathfrak{p}_{x'}$.
		
		\item A point $x \in \Spec R$ is maximal if and only if $\mathfrak{p}_x$ is a minimal prime ideal.
	\end{enumerate}
\end{proposition}

\begin{exercise}
	Prove the proposition.
\end{exercise}

\begin{corollary}\label{chap2:closed-irreducible-subset-of-affine-has-unique-generic-point}
	Let $R$ be a ring. Each closed irreducible subset $Z \subseteq \Spec R$ contains a unique generic point $\eta$ such that $Z = \overline{\Set{\eta}}$.
\end{corollary}
\begin{proof}
	Combine \cref{chap2:correspondence-closed-irreducible-minimal} and \cref{chap2:spacial-points-of-rings}.
\end{proof}

\begin{example}[Spectrum of a PID]
	Let $R$ be a principal ideal domain. The zero ideal $(0)$ is prime and furthermore is a generic point of $\Spec R$. Clearly this is the unique generic point of $\Spec R$.
	
	The other prime ideals of $R$ are ideals generated by prime elements. Since prime elements and irreducible elements coincide in a PID, it follows that every point of $\Spec R$ distinct from the generic point is closed.
\end{example}

\section{Spec as a Functor}

We have seen that a ring $R$ gives rise to the topological space $\Spec R$. In addition to this, a ring homomorphism $\varphi \colon R \to S$ gives rise to a continuous map between the underlying topological space $\Spec S \to \Spec R$, as we shall see in the following lemma.

\begin{lemma}\label{chap2:spec-map-is-continuous}
	Let $\varphi \colon R \to S$ be a ring homomorphism. Then there is a continuous map $\Spec \varphi \colon \Spec S \to \Spec R$ given by $\mathfrak{q}  \mapsto \varphi^{-1}(\mathfrak{q})$. 
	
	In particular, the identity $(\Spec \varphi)^{-1}(V(\mathfrak{b})) = V(\varphi(\mathfrak{b}))$ holds for all ideals $\mathfrak{b}$ of $R$.
\end{lemma}
\begin{proof}
	It is an easy exercise in commutative algebra to show that $\varphi^{-1}(\mathfrak{q})$ is a prime ideal of $R$ if $\mathfrak{q}$ is a prime ideal of $S$. To prove that $f \coloneqq \Spec \varphi$ is continuous, let $V(\mathfrak{b})$ be a closed subset of $\Spec R$. Then
	\begin{equation*}
		f^{-1}(V(\mathfrak{b})) = \Set{\mathfrak{p} \in \Spec S : \mathfrak{b} \subseteq \varphi^{-1}(\mathfrak{p}) },
	\end{equation*}
	from the definition, as one can check. But $\mathfrak{b} \subseteq \varphi^{-1}(\mathfrak{p})$ if and only if $\varphi(\mathfrak{b}) \subseteq \mathfrak{p}$, and thus $f^{-1}(V(\mathfrak{b})) = V(\varphi(\mathfrak{b}))$ as claimed. Since every preimage under $f$ of a closed set is closed, this shows that $f$ is continuous.
\end{proof}

We left the reader to check that, for any ring $R$ and the identity map $\id_R \colon R \to R$, the map $\Spec(\id_R)$ is the identity function of $\Spec R$. Moreover, we have $\Spec(\psi \circ \varphi) = \Spec\varphi \circ \Spec \psi$ for any ring homomorphisms $\varphi\colon R \to S$ and $\psi \colon S \to T$. Hence $\Spec$ is a contravariant functor from the category of rings to the category of topological spaces.

A continuous map between topological spaces $f \colon X \to Y$ is \inlinedef{dominant} if $f(X)$ is dense in $Y$. For the morphism $\Spec \varphi \colon \Spec S \to \Spec R$ between spectra, this happens precisely when $\ker \varphi$ is contained in the nilradical of $R$.

\begin{corollary}\label{chap2:dominant-iff-kernel-in-nilradical}
	Let $\varphi \colon R \to S$ be a ring homomorphism. Then $\Spec \varphi$ is dominant if and only if $\ker \varphi \subseteq \Nil(R)$.
\end{corollary}
\begin{proof}
	See \cref{chap2:dominant-iff-kernel-in-nilradical-problem}.
\end{proof}

In particular, an injective ring map $\varphi \colon R \to S$ induces a dominant continuous map $\Spec \varphi$ on spectra.

\begin{proposition}
	Let $\varphi \colon R \to S$ be a surjective homomorphism of rings with kernel $\mathfrak{a}$. Then $\Spec \varphi$ defines a homeomorphism from $\Spec S$ to the set $V(\mathfrak{a})$ of $\Spec R$.
\end{proposition}

\begin{proof}
	By the first isomorphism theorem, $S \cong R/\mathfrak{a}$. It is then clear that $\Spec \varphi$ is a bijection between prime ideals of $S$ and prime ideals of $R$ containing $\mathfrak{a}$. Hence $\Spec \varphi \colon \Spec S \to V(\mathfrak{a})$ is a bijective continuous map. Now let $\mathfrak{b}$ be an ideal of $S$. Then $\varphi^{-1}(\mathfrak{b})$ is an ideal of $R$. The identity in \cref{chap2:spec-map-is-continuous} takes the form 
	\begin{equation*}
		(\Spec \varphi)^{-1}(V(\varphi^{-1}(\mathfrak{b}))) = V(\varphi(\varphi^{-1}(\mathfrak{b}))).
	\end{equation*}
	Since $\varphi$ is surjective, $\varphi(\varphi^{-1}(A)) = A$ for all $A \subseteq S$. Thus we obtain
	\begin{equation*}
		(\Spec \varphi)^{-1}(V(\varphi^{-1}(\mathfrak{b}))) = V(\mathfrak{b}).
	\end{equation*}
	Using the fact that $\mathfrak{a} \subseteq \varphi^{-1}(\mathfrak{b})$, we get $V(\varphi^{-1}(\mathfrak{b})) \subseteq V(\mathfrak{a})$. Together with $\Spec \varphi \colon \Spec S \to V(\mathfrak{a})$ being bijective, we conclude that $V(\varphi^{-1}(\mathfrak{b})) = (\Spec \varphi) (V(\mathfrak{b}))$. Thus $\Spec \varphi$ is a closed map, and so it is a homeomorphism.
\end{proof}

Since every quotient homomorphism is surjective, we obtain the following corollary.

\begin{corollary}[Homeomorphism between $\Spec R/\mathfrak{a}$ and $V(\mathfrak{a})$]
	Let $R$ be a ring, $\mathfrak{a}$ its ideal, and $\pi \colon R \to R/\mathfrak{a}$ the quotient homomorphism. Then $\Spec \pi$ is a homeomorphism from $\Spec R/\mathfrak{a}$ onto the set $V(\mathfrak{a})$.
\end{corollary}

\begin{proposition}\label{chap2:localisations-isomorphic-to-spec}
	Let $R$ be a ring, $S \subseteq R$ a multiplicative subset, and $\psi \colon R \to S^{-1}R$ the canonical homomorphism $r \mapsto r/1$. Then $\Spec \psi$ is a homeomorphism from $\Spec(S^{-1}R)$ onto the set of prime ideals $\mathfrak{p}$ of $R$ with $\mathfrak{p} \cap S = \emptyset$.
\end{proposition}

\begin{proof}
	Let $D(S)$ denote the set of prime ideals $\mathfrak{p}$ of $R$ with $\mathfrak{p} \cap S = \emptyset$. It is a result in commutative algebra which states that $\Spec \psi$ is a bijection from $\Spec(S^{-1}R)$ to $D(S)$. We show that it is closed. Let $J$ be an ideal of $S^{-1}R$. Then it is easily checked that
	\begin{equation*}
		(\Spec \psi)(V(J)) = V(\psi^{-1}(J)) \cap D(S),
	\end{equation*}
	proving that $\Spec \psi$ is a closed map.
\end{proof}

\section{Problems}

\begin{problems}
	\item Prove that, for any ring $R$ and any ideal $I$ of $R$, $V(I) = \emptyset$ if and only if $I = R$.
	
	\item Prove the following theorem from commutative algebra: if $R$ is a ring and $\mathfrak{a}$ is an ideal of $R$, then
	\begin{equation*}
		\sqrt{\mathfrak{a}} = \bigcap_{\substack{\mathfrak{p} \text{ prime} \\ \mathfrak{p} \supseteq \mathfrak{a}}} \mathfrak{p}.
	\end{equation*}
	
	\item Show that if $f \in R$ vanishes at every point of $\Spec R$ (that is, $f(\mathfrak{p}) = 0$ for all $\mathfrak{p} \in \Spec R$), then $f$ is nilpotent.
	
	
	\item Prove that $\Spec R$ is irreducible if and only if the nilradical of $R$ is prime.
	
	\item Show that if $\mathfrak{p} \subseteq R$ is a prime ideal, then $\Set{\mathfrak{p}}$ is the only generic point of $V(\mathfrak{p})$.
	
	\item Show that if $R$ is an integral domain, then $\Spec R$ is irreducible.
	
	\item \label{chap2:property-of-spec-map} Given a ring homomorphism $\varphi \colon R \to S$, let $f$ denote the corresponding continuous map $\Spec \varphi \colon \Spec S \to \Spec R$. Show that 
	\begin{problems}
		\item $f^{-1}(D(r)) = D(\varphi(r))$ for all $r \in R$.
		
		\item $V(\varphi^{-1}(\mathfrak{b})) = \overline{f(V(\mathfrak{b}))}$ for all ideal $\mathfrak{b} \subseteq S$.
	\end{problems}
	
	\item \label{chap2:dominant-iff-kernel-in-nilradical-problem} Prove \cref{chap2:dominant-iff-kernel-in-nilradical}. \textit{[Hint: Use \cref{chap2:property-of-spec-map}].}
	
	\item Let $R$ be a ring and $\mathfrak{p}$ a prime ideal of $R$. 
	\begin{problems}
		\item Prove that $\Spec R_{\mathfrak{p}}$ has a unique closed point.
		
		\item Prove that if $U$ is an open subset which contains the unique closed point of $\Spec R_{\mathfrak{p}}$, then $U = \Spec R_{\mathfrak{p}}$.
	\end{problems}
	
	\item Let $\varphi \colon R \to S$ be a ring homomorphism. If $\Spec S$ is irreducible with a generic point $\eta$, prove that $(\Spec\varphi)(\eta)$ is the generic point of the closure $\overline{\Img\Spec}$.
	
	\item Show that if $R$ is a ring, then $\Spec R$ and $\Spec(R/\Nil{R})$ are homeomorphic.
\end{problems}

